\sect{उन्मीलनीव्रतम्}
\label{sec:padma-unmilani}

\ganapatyadiDhyanam
\padmaPuranam

\dnsub{अथ सप्तत्रिंशोऽध्यायः॥३७}

\uvacha{महादेव उवाच}

\twolineshloka
{अतस्त्वां सम्प्रवक्ष्यामि उन्मीलनीमनुत्तमाम्}
{यस्याः श्रवणमात्रेण जन्मसंसारबन्धनात्}% ॥१॥

\twolineshloka
{पापात्मा मुच्यते पापैः स्वर्गलोके महीयते}
{देवताः पितरश्चैव तस्या गतिमवाप्नुयुः}% ॥२॥

\twolineshloka
{विद्यार्थी लभते विद्यां सर्वकामानवाप्नुयात्}
{तस्या व्रतान्न सन्देहः स्वर्गलोके महीयते}% ॥३॥

\twolineshloka
{स्वस्थानं तत्र वै प्राप्तः शिवलोके महीयते}
{अतस्त्वं कुरुषे राजन् वैष्णवानां तु पूजनम्}% ॥४॥

\twolineshloka
{वैष्णवानां तु ये राजन् सेवा कुर्वन्ति नित्यशः}
{तेषां दण्डं च कुरुषे नो वा तेषां नराधिपः}% ॥५॥

\twolineshloka
{भोजनानन्तरं तेषां भोजनं कुरुते नृप}
{तैरेव पूजितो विष्णुर्यैर्भक्त्या तु प्रपूजितः}% ॥६॥

\twolineshloka
{शालग्रामशिलाभूतां दत्त्वा मूर्धनि प्रत्यहम्}
{त्वं धारयसि भूपाल कण्ठे नित्यं सुभक्तितः}% ॥७॥

\twolineshloka
{धूपशेषं तु वै विष्णोर्भक्त्या भजसि भूपते}
{आरार्तिकं सदा कृत्वा भक्तानां वेदयेर्नृप}% ॥८॥

\twolineshloka
{शङ्खोदकं हरेर्मूर्ध्नि भ्रामयित्वा तु भक्तितः}
{नित्यं बिभर्षि शिरसि शेषं यच्छसि वैष्णवान्}% ॥९॥

\twolineshloka
{नैवेद्यं प्रत्यहं कृत्वा सर्वोपस्करसंयुतम्}
{विष्वक्सेनाय दत्वा वै स्वयं भुनक्षि वाडव}% ॥१०॥

\twolineshloka
{विष्णोर्निवेदितं चान्नं वैष्णवैः सह भुज्यते}
{नित्यं नामसहस्रेण भक्त्या स्तौषि जनार्दनम्}% ॥११॥

\twolineshloka
{दीपार्घदानं वै विष्णोः कुरुषे गीतनर्तनम्}
{श्यामाङ्कुरैः पूजयसे पूज्यन्ते नृपसत्तम}% ॥१२॥

\twolineshloka
{श्यामाङ्कुरैः सदा वत्स पूजनं चाति दुर्लभम्}
{पृथ्वीदानसमं पुण्यं दूर्वया पूजने कृते}% ॥१३॥

\twolineshloka
{अतो वै नास्ति लोकेऽस्मिन् दूर्वायाः सदृशं भुवि}
{तया वै पूजनं कार्यं विष्णुसायुज्यमिच्छता}% ॥१४॥

\twolineshloka
{अतस्त्वं कुरुषे नित्यं पूजनं दूर्वया सह}
{यवाक्षतैर्विशेषेण पूजनं कुरुषे न वा}% ॥१५॥

\twolineshloka
{पक्षे पक्षे नृपश्रेष्ठ विधिवद् द्वादशीव्रतम्}
{यत् कृतं तु महाराज महापापप्रणाशनम्}% ॥१६॥

\twolineshloka
{मोक्षदं सुखदं चैव तथाऽयुष्यप्रदं सदा}
{एतद् विष्णुव्रतं प्रोक्तं वैष्णवानां तु मोक्षदम्}% ॥१७॥

\twolineshloka
{गृहस्थानां तु सुखदं यतीनां मुक्तिदायकम्}
{सर्वरोगादिशमनं पवित्रं कायशोधनम्}% ॥१८॥

\twolineshloka
{व्रतमेतच्च कुरुषे नो वा चैव नराधिप}
{दशमीवेधरहितं कुरुषे जागरान्वितम्}% ॥१९॥

\twolineshloka
{तुलसीपत्रनिकरैर्नित्यं पूजयसे हरिम्}
{गोपीचन्दनजं पत्रं भाले वा नृपसत्तम}% ॥२०॥

\twolineshloka
{धारितं सर्वलोकानां पवित्रीकरणं नृप}
{अतस्त्वं च धारयसे गोपीचन्दनसम्भवम्}% ॥२१॥

\twolineshloka
{ब्रह्महा हेमहारी च मद्यपानी तथैव च}
{अगम्यगो महापापी तथा ह्यनृतभाषितः}% ॥२२॥

\twolineshloka
{ते सर्वे मुक्तिमायान्ति तिलकं धारणादृताः}
{बिभर्षि कण्ठे नित्यं त्वं धात्रीफलसमुद्भवाम्}% ॥२३॥

\twolineshloka
{मालां मुख्यायुतसमां तुलसीपत्रसम्भवाम्}
{शालग्रामशिलायुक्तां द्वारकायां समुद्भवाम्}% ॥२४॥

\twolineshloka
{नित्यं पूजयसे भूप भुक्तिमुक्तिफलप्रदाम्}
{पद्मसंज्ञं पुराणं वै पठसे पुरतो हरेः}% ॥२५॥

\twolineshloka
{चरितं दैत्यराज्यस्य प्रह्लादस्य च भूपते}
{वासरं वासुदेवस्य सवेधं कुर्वतो नरान्}% ॥२६॥

\twolineshloka
{निवारयसि भूपाल शास्त्रं दृष्ट्वा प्रयत्नतः}
{सवेधं वासरं विष्णोर्यस्मिन् राष्ट्रे प्रवर्तते}% ॥२७॥

\threelineshloka
{लिप्यते तेन पापेन राजा भवति नारकी}
{वेधं चतुर्विधं त्यक्त्वा समुपोष्य हरेर्दिनम्}
{कुलकोटिं समुद्धृत्य विष्णुलोके महीयते}% ॥२८॥

\uvacha{गौतम उवाच}

\twolineshloka
{शृणु भूपाल वक्ष्यामि वैष्णवाख्यं महाव्रतम्}
{यं श्रुत्वा पापिनः सर्वे मुक्तिमायान्ति तत्क्षणात्}% ॥२९॥

\twolineshloka
{द्वादशीसम्भवं पुण्यं मया ख्यातं न कस्यचित्}
{वैष्णवोऽसि महाराज भक्तो भागवते नृणाम्}% ॥३०॥

\twolineshloka
{वैष्णवं तु महागुह्यं तद् व्रतं त्वं निशामय}
{उन्मीलनी नाम पुरा भक्त्या मे माधवेन तु}% ॥३१॥

\twolineshloka
{कथिता सुप्रसन्नेन तां ते भूप वदाम्यहम्}
{एकादशी अहोरात्रं प्रभाते घटिका भवेत्}% ॥३२॥

\twolineshloka
{उन्मीलनी तु सा ज्ञेया विशेषेण हरिप्रिया}
{त्रैलोक्ये यानि तीर्थानि पुण्यान्यायतनानि च}% ॥३३॥

\twolineshloka
{कोट्यंशे नैव तुल्यानि मखा वेदास्तपांसि च}
{उन्मीलनी समं किञ्चिन्न भूतं न भविष्यति}% ॥३४॥

\twolineshloka
{प्रयागो न कुरुक्षेत्रं न काशी न च पुष्करः}
{शैलो हिमाचलो नैव न मेरुर्गन्धमादनः}% ॥३५॥

\twolineshloka
{शैलो न नीलनिषधो न विन्ध्यो नैव नैमिषम्}
{गोदावरी न कावेरी चन्द्रभागा न वेदिका}% ॥३६॥

\twolineshloka
{न तापी न पयोष्णी च न क्षिप्रा नैव चन्दना}
{चर्मण्वती च सरयूश्चन्द्रभागा न गण्डिका}% ॥३७॥

\twolineshloka
{गोमती च विपाशा च शोणाख्यश्च महानदः}
{किमत्र बहुनोक्तेन भूयो भूयो नराधिप}% ॥३८॥

\twolineshloka
{उन्मीलनीसमं किञ्चिन्न देवः केशवात् परः}
{उन्मीलनीमनुप्राप्य यैः कृतं केशवार्चनम्}% ॥३९॥

\twolineshloka
{पापचक्रसमूहस्य राशयः पतिताः क्षणात्}
{यस्मिन् मासे महीपाल तिथिरुन्मीलनी भवेत्}% ॥४०॥

\twolineshloka
{तन्मासनाम्ना गोविन्दः पूजनीयः प्रयत्नतः}
{जातरूपमयः कार्यं मासनाम्ना तु माधवः}% ॥४१॥

\twolineshloka
{स्वशक्त्या विश्वरूपस्तु श्रद्धाभक्तिसमन्वितः}
{पवित्रोदकसंयुक्तं पञ्चरत्नसमन्वितम्}% ॥४२॥

\twolineshloka
{गन्धपुष्पाक्षतैर्युक्तं कुम्भं स्रग्दामभूषितम्}
{पात्रं च सोदकं कार्यं गोधूमैश्चापि पूरितम्}% ॥४३॥

\twolineshloka
{नानारत्नैश्च संयुक्तं नानागन्धैः प्रपूजितम्}
{मल्लिकामोदसंयुक्तं जातीपुष्पैः प्रपूजितम्}% ॥४४॥

\twolineshloka
{श्वेताख्यैस्तन्दुलैश्चैव पूरणीयः प्रयत्नतः}
{प्रदद्याद् वस्त्रयुग्मं तु उपवीतं तु सोत्तरम्}% ॥४५॥

\twolineshloka
{उपानहौ तु राजर्षे आतपत्रं शिरोपरि}
{भोजनं जलपात्रं च सप्तधान्यं तिलैः सह}% ॥४६॥

\twolineshloka
{रूप्यं चैव तु कार्पासं पायसं मुद्रिका हरेः}
{धेनुर्वाऽत्र तु दातव्या वत्सालङ्कारसंयुता}% ॥४७॥

\twolineshloka
{स्वर्णशृङ्गी रौप्यखुरी ताम्रपृष्ठी तथैव च}
{कांस्यदोहीं रत्नपुच्छीं दद्याद्वै गुरवे तदा}% ॥४८॥

\twolineshloka
{शय्यां सोपस्करां दद्यात् साधवे भक्तिपूर्वकम्}
{धूपं दीपं तु नैवेद्यं फलपत्रं निवेदयेत्}% ॥४९॥

\twolineshloka
{पूजनीयो महाभक्तैर्मन्त्रैरेभिस्तु केशवः}
{तुलसीपत्रसंयुक्तैः पुष्पैः कालोद्भवैस्तथा}% ॥५०॥

\twolineshloka
{मासनाम्नैव चरणौ जानुनी विष्णुरूपिणे}
{गुह्ये तु भूतपतये कटिं वै पीतवाससे}% ॥५१॥

\twolineshloka
{ब्रह्ममूर्तिभृते नाभिमुदरं विश्वयोनये}
{हृदयं ज्ञानगम्याय कण्ठं वैकुण्ठमूर्तये}% ॥५२॥

\twolineshloka
{ऊर्ध्वगाय ललाटे तु बाहौ दक्षान्तकारिणे}
{उत्तमाङ्गं सुरेशाय सर्वाङ्गं सर्वमूर्तये}% ॥५३॥

\twolineshloka
{स्वनाम्ना चाऽऽयुधादीनि पूजनीयानि भक्तितः}
{अर्घ्यदानं प्रकर्तव्यं नालिकेरादिभिः समम्}% ॥५४॥

\twolineshloka
{शङ्खोपरि जले कृत्वा गन्धपुष्पाक्षतान्वितम्}
{सूत्रेणावेष्टितं कृत्वा दद्यादर्घ्यं विधानतः}% ॥५५॥

\twolineshloka
{देवदेव महादेव श्रीकेशव जनार्दन}
{सुब्रह्मण्य नमस्तेऽस्तु पुण्यराशिविवर्धन}% ॥५६॥

\twolineshloka
{शोकमोहमहापापान्मामुद्धर भवार्णवात्}
{सुकृतं न कृतं किञ्चिज्जन्मकोटिशतैरपि}% ॥५७॥

\twolineshloka
{तथाऽपि मां महास्वामिन् मामुद्धर भवार्णवात्}
{व्रतेनानेन देवेश ये चान्ये मम पूर्वजाः}% ॥५८॥

\twolineshloka
{वियोनिं च गताश्चान्ये पापमृत्युवशं गताः}
{ये भविष्यन्ति येऽतीताः प्रेतलोकान् समुद्धर}% ॥५९॥

\twolineshloka
{आर्तस्य मम दीनस्य भक्तिरव्यभिचारिणी}
{दत्तमर्घ्यं मया तुभ्यं भक्त्या गृह्ण गदाधर}% ॥६०॥

\twolineshloka
{दत्त्वाऽर्घ्यं धूपदीपाद्यैर्नैवेद्यैर्विष्णुसम्भवैः}
{स्तोत्रैर्नीराजनैर्गीतैर्भृत्यैः सन्तोषयेद्धरिम्}% ॥६१॥

\twolineshloka
{वस्त्रैर्दानैश्च गोदानैर्भोजनैस्तोषयेद् गुरुम्}
{तथा तथा विधातव्यं गुरुर्वै प्रीतिमाप्नुयात्}% ॥६२॥

\twolineshloka
{लोकानां तारणार्थाय धात्रा सृष्टो गुरुर्यतः}
{अतो वै गुरुपूजा च कर्तव्या वै प्रयत्नतः}% ॥६३॥

\twolineshloka
{अहितं यो नाशयति स्वहितं दर्शयेत् सदा}
{स गुरुः स च विज्ञेयः सर्वधर्मार्थकोविदः}% ॥६४॥

\twolineshloka
{अकुर्वन् वित्तशाठ्यं तु गुरवे तं निवेदयेत्}
{गुरोर्निवेदिते भूप परिपूर्णं भवेद् व्रतम्}% ॥६५॥

\twolineshloka
{कृत्वा दिनकृतं कर्म भोजनं ब्राह्मणैः सह}
{कर्तव्यं नृपशार्दूल दिनं नेयं कथानकैः}% ॥६६॥

\twolineshloka
{अनेन विधिना यस्तु कुर्यादुन्मीलनीव्रतम्}
{कल्पकोटिसहस्राणि वसते विष्णुसन्निधौ}% ॥६७॥

॥इति श्रीपाद्मे महापुराणे पञ्चपञ्चाशत्साहस्र्यां संहितायामुत्तरखण्डे उमापतिनारदसंवादान्तर्गतकृष्णयुधिष्ठिरसंवादे उन्मीलनीव्रतं नाम सप्तत्रिंशोऽध्यायः॥३७॥

\genMangalaShloka

\hyperref[sec:ekadashi_mahatmyam_padma_puranam]{\closesub}
\clearpage

